\documentclass[10pt,twoside,openany]{memoir}
%\usepackage{mlmodern}
%\usepackage{tgpagella} % text only
%\usepackage{mathpazo}  % math & text
\usepackage[T1]{fontenc}
\usepackage[hidelinks]{hyperref}
\usepackage{amsmath}
\usepackage[fixamsmath]{mathtools}  % Extension to amsmath
\usepackage{amsthm}
\usepackage{amssymb}

\usepackage{newpxtext}
\usepackage{eulerpx}
\usepackage{eucal}
\usepackage{datetime}
    \newdateformat{specialdate}{\THEYEAR\ \monthname\ \THEDAY}
\usepackage[margin=1.5in]{geometry}
\usepackage{fancyhdr}
    \pagestyle{fancy}
    \renewcommand{\headrulewidth}{0pt}
    \cfoot{\scriptsize \thepage}
\usepackage{thmtools}
    \declaretheoremstyle[
        spaceabove=10pt,
        spacebelow=10pt,
        headfont=\normalfont\bfseries,
        notefont=\mdseries, notebraces={(}{)},
        bodyfont=\normalfont,
        postheadspace=0.5em
        %qed=\qedsymbol
        ]{defs}

    \declaretheoremstyle[ 
        spaceabove=10pt, % space above the theorem
        spacebelow=10pt,
        headfont=\normalfont\bfseries,
        bodyfont=\normalfont\itshape,
        postheadspace=0.5em
        ]{thmstyle}
    
    \declaretheorem[
        style=thmstyle,
        numberwithin=section
    ]{theorem}

    \declaretheorem[
        style=thmstyle,
        sibling=theorem,
    ]{proposition}

    \declaretheorem[
        style=thmstyle,
        sibling=theorem,
    ]{lemma}

    \declaretheorem[
        style=thmstyle,
        sibling=theorem,
    ]{corollary}

    \declaretheorem[
        numberwithin=section,
        style=defs,
    ]{example}

    \declaretheorem[
        numberwithin=section,
        style=defs,
    ]{definition}

    \declaretheorem[
        style=defs,
        sibling=theorem,
        numberwithin=section,
    ]{exercise}

    \declaretheorem[
        style=thmstyle
        numbered=unless unique
    ]{axiom}

    \declaretheorem[numbered=unless unique, style=defs]{problem}
    \declaretheorem[numberwithin=section,style=defs]{note}
    \declaretheorem[numbered=no,style=defs]{question}
    \declaretheorem[numbered=no,style=defs]{recall}
    \declaretheorem[numbered=no,style=remark]{answer}
    \declaretheorem[numbered=no,style=remark]{solution}
    \declaretheorem[numbered=no,style=defs]{remark}
\usepackage{enumitem}
\usepackage[utf8x]{inputenc}
\usepackage{tikz}
\usepackage{tikz-cd}
\usetikzlibrary{patterns}
\usetikzlibrary{positioning,arrows.meta}
\linespread{1.00}
%%%%%%%%%%%%%%%%%%%%%%%%%%%%%%%%%%%%%%%%%%%%%%%%%%%%%%%%%%%%%
%%%%%%%%%%%%%%%%%%%%%%%%%%%%%%%%%%%%%%%%%%%%%%%%%%%%%%%%%%%%%
% ============================================================
% makros.tex — reorganized, full restored version
% ============================================================

% ============================================================
% Sha / Cyrillic support notes
% ============================================================

%to make the correct symbol for Sha
%\newcommand\cyr{%
%\renewcommand\rmdefault{wncyr}%
%\renewcommand\sfdefault{wncyss}%
%\renewcommand\encodingdefault{OT2}%
%\normalfont \selectfont} \DeclareTextFontCommand{\textcyr}{\cyr}

% ============================================================
% Math operators
% ============================================================

\DeclareMathOperator{\ab}{ab}
\DeclareMathOperator{\ad}{ad}
\DeclareMathOperator{\adj}{adj}
\DeclareMathOperator{\alg}{alg}
\DeclareMathOperator{\Alt}{Alt}
\DeclareMathOperator{\Ann}{Ann}
\DeclareMathOperator{\arith}{arith}
\DeclareMathOperator{\Aut}{Aut}
\DeclareMathOperator{\Be}{B}
\DeclareMathOperator{\Bd}{Bd}
\DeclareMathOperator{\card}{card}
\DeclareMathOperator{\Char}{char}
\DeclareMathOperator{\csp}{csp}
\DeclareMathOperator{\codim}{codim}
\DeclareMathOperator{\coker}{coker}
\DeclareMathOperator{\coh}{H}
\DeclareMathOperator{\compl}{compl}
\DeclareMathOperator{\conj}{conj}
\DeclareMathOperator{\cont}{cont}
\DeclareMathOperator{\Cov}{Cov}
\DeclareMathOperator{\crys}{crys}
\DeclareMathOperator{\Crys}{Crys}
\DeclareMathOperator{\cusp}{cusp}
\DeclareMathOperator{\diag}{diag}
\DeclareMathOperator{\diam}{diam}
\DeclareMathOperator{\Dom}{Dom}
\DeclareMathOperator{\disc}{disc}
\DeclareMathOperator{\dist}{dist}
\DeclareMathOperator{\Div}{div}
\DeclareMathOperator{\dR}{dR}
\DeclareMathOperator{\Eis}{Eis}
\DeclareMathOperator{\End}{End}
\DeclareMathOperator{\ev}{ev}
\DeclareMathOperator{\eval}{eval}
\DeclareMathOperator{\Eq}{Eq}
\DeclareMathOperator{\Ext}{Ext}
\DeclareMathOperator{\Fil}{Fil}
\DeclareMathOperator{\Fitt}{Fitt}
\DeclareMathOperator{\Frob}{Frob}
\DeclareMathOperator{\G}{G}
\DeclareMathOperator{\Gal}{Gal}
\DeclareMathOperator{\GL}{GL}
\DeclareMathOperator{\Gr}{Gr}
\DeclareMathOperator{\Graph}{Graph}
\DeclareMathOperator{\GSp}{GSp}
\DeclareMathOperator{\GUn}{GU}
\DeclareMathOperator{\Hom}{Hom}
\DeclareMathOperator{\id}{id}
\DeclareMathOperator{\Id}{Id}
\DeclareMathOperator{\Ik}{Ik}
\DeclareMathOperator{\IM}{Im}
\DeclareMathOperator{\Image}{im}
\DeclareMathOperator{\Ind}{Ind}
\DeclareMathOperator{\Inf}{inf}
\DeclareMathOperator{\ins}{ins}
\DeclareMathOperator{\inv}{inv}
\DeclareMathOperator{\Isom}{Isom}
\DeclareMathOperator{\J}{J}
\DeclareMathOperator{\Jac}{Jac}
\DeclareMathOperator{\lcm}{lcm}
\DeclareMathOperator{\length}{length}
\DeclareMathOperator*{\limit}{limit}
\DeclareMathOperator{\Log}{Log}
\DeclareMathOperator{\M}{M}
\DeclareMathOperator{\Mat}{Mat}
\DeclareMathOperator{\N}{N}
\DeclareMathOperator{\Nm}{Nm}
\DeclareMathOperator{\NIk}{N-Ik}
\DeclareMathOperator{\NSK}{N-SK}
\DeclareMathOperator{\new}{new}
\DeclareMathOperator{\obj}{obj}
\DeclareMathOperator{\old}{old}
\DeclareMathOperator{\ord}{ord}
\DeclareMathOperator{\Or}{O}
\DeclareMathOperator{\op}{op}
\DeclareMathOperator{\out}{out}
\DeclareMathOperator{\PGL}{PGL}
\DeclareMathOperator{\PGSp}{PGSp}
\DeclareMathOperator{\rank}{rank}
\DeclareMathOperator{\Ran}{Ran}
\DeclareMathOperator{\rel}{Rel}
\DeclareMathOperator{\Real}{Re}
\DeclareMathOperator{\RES}{res}
\DeclareMathOperator{\Res}{Res}
%\DeclareMathOperator{\Sha}{\textcyr{Sh}}
\DeclareMathOperator{\Sel}{Sel}
\DeclareMathOperator{\semi}{ss}
\DeclareMathOperator{\sgn}{sign}
\DeclareMathOperator{\SK}{SK}
\DeclareMathOperator{\SL}{SL}
\DeclareMathOperator{\SO}{SO}
\DeclareMathOperator{\Sp}{Sp}
\DeclareMathOperator{\Span}{span}
\DeclareMathOperator{\Spec}{Spec}
\DeclareMathOperator{\spin}{spin}
\DeclareMathOperator{\st}{st}
\DeclareMathOperator{\St}{St}
\DeclareMathOperator{\SUn}{SU}
\DeclareMathOperator{\supp}{supp}
\DeclareMathOperator{\Sup}{sup}
\DeclareMathOperator{\Sym}{Sym}
\DeclareMathOperator{\Tam}{Tam}
\DeclareMathOperator{\tors}{tors}
\DeclareMathOperator{\tr}{tr}
\DeclareMathOperator{\Tr}{Tr}
\DeclareMathOperator{\un}{un}
\DeclareMathOperator{\Un}{U}
\DeclareMathOperator{\val}{val}
\DeclareMathOperator{\vol}{vol}

% ============================================================
% General aliases and short macros
% ============================================================

\newcommand{\absgal}{\G_{\bbQ}}
\newcommand{\Loc}{L_{\operatorname{loc}}}
\renewcommand{\k}{\kappa}
\newcommand{\Ff}{F_{f}}
%\newcommand{\ts}{\,^{t}\!}
\newcommand{\lat}{\mathscr{L}}
\newcommand{\ov}[1]{\overline{#1}}
\newcommand{\up}{\upsilon}
\newcommand{\e}{\epsilon}
\newcommand{\tac}{\textasteriskcentered}

%mahesh macros
\newcommand{\tm}{\textrm}

\newcommand{\bone}{\mathbf{1}}
\newcommand{\sign}{\mathrm{sign}}
\newcommand{\eps}{\varepsilon}
\newcommand{\textui}[1]{\underline{\textit{#1}}}

%\newcommand{\ov}{\overline}
%\newcommand{\un}{\underline}
\newcommand{\fin}{\mathrm{fin}}
\newcommand{\nl}{\newline \mbox{}}
\newcommand{\h}[1]{\hspace{#1pt}}
\newcommand{\mtext}[1]{\hspace{6pt}\text{#1}\hspace{6pt}}
\newcommand{\lnorm}{\left\lVert}
\newcommand{\rnorm}{\right\rVert}
\newcommand{\ds}{\displaystyle}
\newcommand{\ts}{\textstyle}
\newcommand{\sfrac}[2]{{}^{#1}\mskip -5mu/\mskip -3mu_{#2}}

% ============================================================
% Category-style operators and contoured variants
% ============================================================

\DeclareMathOperator{\Sets}{S \mkern1.04mu e \mkern1.04mu t \mkern1.04mu s}
    \newcommand{\cSets}{\scalebox{1.02}{\contour{black}{$\Sets$}}}
    
\DeclareMathOperator{\Groups}{G \mkern1.04mu r \mkern1.04mu o \mkern1.04mu u \mkern1.04mu p \mkern1.04mu s}
    \newcommand{\cGroups}{\scalebox{1.02}{\contour{black}{$\Groups$}}}

\DeclareMathOperator{\TTop}{T \mkern1.04mu o \mkern1.04mu p}
    \newcommand{\cTop}{\scalebox{1.02}{\contour{black}{$\TTop$}}}

\DeclareMathOperator{\Htp}{H \mkern1.04mu t \mkern1.04mu p}
    \newcommand{\cHtp}{\scalebox{1.02}{\contour{black}{$\Htp$}}}

\DeclareMathOperator{\Mod}{M \mkern1.04mu o \mkern1.04mu d}
    \newcommand{\cMod}{\scalebox{1.02}{\contour{black}{$\Mod$}}}

\DeclareMathOperator{\Ab}{A \mkern1.04mu b}
    \newcommand{\cAb}{\scalebox{1.02}{\contour{black}{$\Ab$}}}

\DeclareMathOperator{\Rings}{R \mkern1.04mu i \mkern1.04mu n \mkern1.04mu g \mkern1.04mu s}
    \newcommand{\cRings}{\scalebox{1.02}{\contour{black}{$\Rings$}}}

\DeclareMathOperator{\ComRings}{C \mkern1.04mu o \mkern1.04mu m \mkern1.04mu R \mkern1.04mu i \mkern1.04mu n \mkern1.04mu g \mkern1.04mu s}
    \newcommand{\cComRings}{\scalebox{1.05}{\contour{black}{$\ComRings$}}}

\DeclareMathOperator{\hHom}{H \mkern1.04mu o \mkern1.04mu m}
    \newcommand{\cHom}{\scalebox{1.02}{\contour{black}{$\hHom$}}}

% ============================================================
% Mathcal
% ============================================================

\newcommand{\cA}{\mathcal{A}}
\newcommand{\cB}{\mathcal{B}}
\newcommand{\cC}{\mathcal{C}}
\newcommand{\cD}{\mathcal{D}}
\newcommand{\cE}{\mathcal{E}}
\newcommand{\cF}{\mathcal{F}}
\newcommand{\cG}{\mathcal{G}}
\newcommand{\cH}{\mathcal{H}}
\newcommand{\cI}{\mathcal{I}}
\newcommand{\cJ}{\mathcal{J}}
\newcommand{\cK}{\mathcal{K}}
\newcommand{\cL}{\mathcal{L}}
\newcommand{\cM}{\mathcal{M}}
\newcommand{\cN}{\mathcal{N}}
\newcommand{\cO}{\mathcal{O}}
\newcommand{\cP}{\mathcal{P}}
\newcommand{\cQ}{\mathcal{Q}}
\newcommand{\cR}{\mathcal{R}}
\newcommand{\cS}{\mathcal{S}}
\newcommand{\cT}{\mathcal{T}}
\newcommand{\cU}{\mathcal{U}}
\newcommand{\cV}{\mathcal{V}}
\newcommand{\cW}{\mathcal{W}}
\newcommand{\cX}{\mathcal{X}}
\newcommand{\cY}{\mathcal{Y}}
\newcommand{\cZ}{\mathcal{Z}}

% ============================================================
% Mathscrl% 
%============================================================

\newcommand{\scra}{\mathscr{a}}
\newcommand{\scrA}{\mathscr{A}}
\newcommand{\scrb}{\mathscr{b}}
\newcommand{\scrB}{\mathscr{B}}
\newcommand{\scrc}{\mathscr{c}}
\newcommand{\scrC}{\mathscr{C}}
\newcommand{\scrd}{\mathscr{d}}
\newcommand{\scrD}{\mathscr{D}}
\newcommand{\scre}{\mathscr{e}}
\newcommand{\scrE}{\mathscr{E}}
\newcommand{\scrf}{\mathscr{f}}
\newcommand{\scrF}{\mathscr{F}}
\newcommand{\scrg}{\mathscr{g}}
\newcommand{\scrG}{\mathscr{G}}
\newcommand{\scrh}{\mathscr{h}}
\newcommand{\scrH}{\mathscr{H}}
\newcommand{\scri}{\mathscr{i}}
\newcommand{\scrI}{\mathscr{I}}
\newcommand{\scrj}{\mathscr{j}}
\newcommand{\scrJ}{\mathscr{J}}
\newcommand{\scrk}{\mathscr{k}}
\newcommand{\scrK}{\mathscr{K}}
\newcommand{\scrl}{\mathscr{l}}
\newcommand{\scrL}{\mathscr{L}}
\newcommand{\scrm}{\mathscr{m}}
\newcommand{\scrM}{\mathscr{M}}
\newcommand{\scrn}{\mathscr{n}}
\newcommand{\scrN}{\mathscr{N}}
\newcommand{\scro}{\mathscr{o}}
\newcommand{\scrO}{\mathscr{O}}
\newcommand{\scrp}{\mathscr{p}}
\newcommand{\scrP}{\mathscr{P}}
\newcommand{\scrq}{\mathscr{q}}
\newcommand{\scrQ}{\mathscr{Q}}
\newcommand{\scrr}{\mathscr{r}}
\newcommand{\scrR}{\mathscr{R}}
\newcommand{\scrs}{\mathscr{s}}
\newcommand{\scrS}{\mathscr{S}}
\newcommand{\scrt}{\mathscr{t}}
\newcommand{\scrT}{\mathscr{T}}
\newcommand{\scru}{\mathscr{u}}
\newcommand{\scrU}{\mathscr{U}}
\newcommand{\scrv}{\mathscr{v}}
\newcommand{\scrV}{\mathscr{V}}
\newcommand{\scrw}{\mathscr{w}}
\newcommand{\scrW}{\mathscr{W}}
\newcommand{\scrx}{\mathscr{x}}
\newcommand{\scrX}{\mathscr{X}}
\newcommand{\scry}{\mathscr{y}}
\newcommand{\scrY}{\mathscr{Y}}
\newcommand{\scrz}{\mathscr{z}}
\newcommand{\scrZ}{\mathscr{Z}}

% ============================================================
% Mathfrak (missing \fi)
% ============================================================

\newcommand{\fa}{\mathfrak{a}}
\newcommand{\fA}{\mathfrak{A}}
\newcommand{\fb}{\mathfrak{b}}
\newcommand{\fB}{\mathfrak{B}}
\newcommand{\fc}{\mathfrak{c}}
\newcommand{\fC}{\mathfrak{C}}
\newcommand{\fd}{\mathfrak{d}}
\newcommand{\fD}{\mathfrak{D}}
\newcommand{\fe}{\mathfrak{e}}
\newcommand{\fE}{\mathfrak{E}}
\newcommand{\ff}{\mathfrak{f}}
\newcommand{\fF}{\mathfrak{F}}
\newcommand{\fg}{\mathfrak{g}}
\newcommand{\fG}{\mathfrak{G}}
\newcommand{\fh}{\mathfrak{h}}
\newcommand{\fH}{\mathfrak{H}}
\newcommand{\fI}{\mathfrak{I}}
\newcommand{\fj}{\mathfrak{j}}
\newcommand{\fJ}{\mathfrak{J}}
\newcommand{\fk}{\mathfrak{k}}
\newcommand{\fK}{\mathfrak{K}}
\newcommand{\fl}{\mathfrak{l}}
\newcommand{\fL}{\mathfrak{L}}
\newcommand{\fm}{\mathfrak{m}}
\newcommand{\fM}{\mathfrak{M}}
\newcommand{\fn}{\mathfrak{n}}
\newcommand{\fN}{\mathfrak{N}}
\newcommand{\fo}{\mathfrak{o}}
\newcommand{\fO}{\mathfrak{O}}
\newcommand{\fp}{\mathfrak{p}}
\newcommand{\fP}{\mathfrak{P}}
\newcommand{\fq}{\mathfrak{q}}
\newcommand{\fQ}{\mathfrak{Q}}
\newcommand{\fr}{\mathfrak{r}}
\newcommand{\fR}{\mathfrak{R}}
\newcommand{\fs}{\mathfrak{s}}
\newcommand{\fS}{\mathfrak{S}}
\newcommand{\ft}{\mathfrak{t}}
\newcommand{\fT}{\mathfrak{T}}
\newcommand{\fu}{\mathfrak{u}}
\newcommand{\fU}{\mathfrak{U}}
\newcommand{\fv}{\mathfrak{v}}
\newcommand{\fV}{\mathfrak{V}}
\newcommand{\fw}{\mathfrak{w}}
\newcommand{\fW}{\mathfrak{W}}
\newcommand{\fx}{\mathfrak{x}}
\newcommand{\fX}{\mathfrak{X}}
\newcommand{\fy}{\mathfrak{y}}
\newcommand{\fY}{\mathfrak{Y}}
\newcommand{\fz}{\mathfrak{z}}
\newcommand{\fZ}{\mathfrak{Z}}

% ============================================================
% Mathbf
% ============================================================

\newcommand{\bfA}{\mathbf{A}}
\newcommand{\bfB}{\mathbf{B}}
\newcommand{\bfC}{\mathbf{C}}
\newcommand{\bfD}{\mathbf{D}}
\newcommand{\bfE}{\mathbf{E}}
\newcommand{\bfF}{\mathbf{F}}
\newcommand{\bfG}{\mathbf{G}}
\newcommand{\bfH}{\mathbf{H}}
\newcommand{\bfI}{\mathbf{I}}
\newcommand{\bfJ}{\mathbf{J}}
\newcommand{\bfK}{\mathbf{K}}
\newcommand{\bfL}{\mathbf{L}}
\newcommand{\bfM}{\mathbf{M}}
\newcommand{\bfN}{\mathbf{N}}
\newcommand{\bfO}{\mathbf{O}}
\newcommand{\bfP}{\mathbf{P}}
\newcommand{\bfQ}{\mathbf{Q}}
\newcommand{\bfR}{\mathbf{R}}
\newcommand{\bfS}{\mathbf{S}}
\newcommand{\bfT}{\mathbf{T}}
\newcommand{\bfU}{\mathbf{U}}
\newcommand{\bfV}{\mathbf{V}}
\newcommand{\bfW}{\mathbf{W}}
\newcommand{\bfX}{\mathbf{X}}
\newcommand{\bfY}{\mathbf{Y}}
\newcommand{\bfZ}{\mathbf{Z}}

\newcommand{\bfa}{\mathbf{a}}
\newcommand{\bfb}{\mathbf{b}}
\newcommand{\bfc}{\mathbf{c}}
\newcommand{\bfd}{\mathbf{d}}
\newcommand{\bfe}{\mathbf{e}}
\newcommand{\bff}{\mathbf{f}}
\newcommand{\bfg}{\mathbf{g}}
\newcommand{\bfh}{\mathbf{h}}
\newcommand{\bfi}{\mathbf{i}}
\newcommand{\bfj}{\mathbf{j}}
\newcommand{\bfk}{\mathbf{k}}
\newcommand{\bfl}{\mathbf{l}}
\newcommand{\bfm}{\mathbf{m}}
\newcommand{\bfn}{\mathbf{n}}
\newcommand{\bfo}{\mathbf{o}}
\newcommand{\bfp}{\mathbf{p}}
\newcommand{\bfq}{\mathbf{q}}
\newcommand{\bfr}{\mathbf{r}}
\newcommand{\bfs}{\mathbf{s}}
\newcommand{\bft}{\mathbf{t}}
\newcommand{\bfu}{\mathbf{u}}
\newcommand{\bfv}{\mathbf{v}}
\newcommand{\bfw}{\mathbf{w}}
\newcommand{\bfx}{\mathbf{x}}
\newcommand{\bfy}{\mathbf{y}}
\newcommand{\bfz}{\mathbf{z}}

% ============================================================
% Blackboard bold
% ============================================================

\newcommand{\bbA}{\mathbb{A}}
\newcommand{\bbB}{\mathbb{B}}
\newcommand{\bbC}{\mathbb{C}}
\newcommand{\bbD}{\mathbb{D}}
\newcommand{\bbE}{\mathbb{E}}
\newcommand{\bbF}{\mathbb{F}}
\newcommand{\bbG}{\mathbb{G}}
\newcommand{\bbH}{\mathbb{H}}
\newcommand{\bbI}{\mathbb{I}}
\newcommand{\bbJ}{\mathbb{J}}
\newcommand{\bbK}{\mathbb{K}}
\newcommand{\bbL}{\mathbb{L}}
\newcommand{\bbM}{\mathbb{M}}
\newcommand{\bbN}{\mathbb{N}}
\newcommand{\bbO}{\mathbb{O}}
\newcommand{\bbP}{\mathbb{P}}
\newcommand{\bbQ}{\mathbb{Q}}
\newcommand{\bbR}{\mathbb{R}}
\newcommand{\bbS}{\mathbb{S}}
\newcommand{\bbT}{\mathbb{T}}
\newcommand{\bbU}{\mathbb{U}}
\newcommand{\bbV}{\mathbb{V}}
\newcommand{\bbW}{\mathbb{W}}
\newcommand{\bbX}{\mathbb{X}}
\newcommand{\bbY}{\mathbb{Y}}
\newcommand{\bbZ}{\mathbb{Z}}

\newcommand{\Bmu}{\mbox{$\raisebox{-0.59ex}
  {$l$}\hspace{-0.18em}\mu\hspace{-0.88em}\raisebox{-0.98ex}{\scalebox{2}
  {$\color{white}.$}}\hspace{-0.416em}\raisebox{+0.88ex}
  {$\color{white}.$}\hspace{0.46em}$}{}}  %need graphicx and xcolor. this produces blackboard bold mu 

% ============================================================
% Greek-letter aliases
% ============================================================

\newcommand{\jota}{\jmath}

% ============================================================
% Matrices
% ============================================================

\newcommand{\bmat}{\left( \begin{matrix}}
\newcommand{\emat}{\end{matrix} \right)}

\newcommand{\bbmat}{\left[ \begin{matrix}}
\newcommand{\ebmat}{\end{matrix} \right]}

\newcommand{\pmat}{\left( \begin{smallmatrix}}
\newcommand{\epmat}{\end{smallmatrix} \right)}

\newcommand{\mat}[4]{\begin{pmatrix}{#1}&{#2}\\{#3}&{#4}\end{pmatrix}}

% ============================================================
% Residues and averaged integrals
% ============================================================

\newcommand{\res}[1]{\underset{#1}{\RES}\,}

\makeatletter
\newcommand*{\mint}[1]{%
  % #1: overlay symbol
  \mint@l{#1}{}%
}
\newcommand*{\mint@l}[2]{%
  % #1: overlay symbol
  % #2: limits
  \@ifnextchar\limits{%
    \mint@l{#1}%
  }{%
    \@ifnextchar\nolimits{%
      \mint@l{#1}%
    }{%
      \@ifnextchar\displaylimits{%
        \mint@l{#1}%
      }{%
        \mint@s{#2}{#1}%
      }%
    }%
  }%
}
\newcommand*{\mint@s}[2]{%
  % #1: limits
  % #2: overlay symbol
  \@ifnextchar_{%
    \mint@sub{#1}{#2}%
  }{%
    \@ifnextchar^{%
      \mint@sup{#1}{#2}%
    }{%
      \mint@{#1}{#2}{}{}%
    }%
  }%
}
\def\mint@sub#1#2_#3{%
  \@ifnextchar^{%
    \mint@sub@sup{#1}{#2}{#3}%
  }{%
    \mint@{#1}{#2}{#3}{}%
  }%
}
\def\mint@sup#1#2^#3{%
  \@ifnextchar_{%
    \mint@sup@sub{#1}{#2}{#3}%
  }{%
    \mint@{#1}{#2}{}{#3}%
  }%
}
\def\mint@sub@sup#1#2#3^#4{%
  \mint@{#1}{#2}{#3}{#4}%
}
\def\mint@sup@sub#1#2#3_#4{%
  \mint@{#1}{#2}{#4}{#3}%
}
\newcommand*{\mint@}[4]{%
  % #1: \limits, \nolimits, \displaylimits
  % #2: overlay symbol: -, =, ...
  % #3: subscript
  % #4: superscript
  \mathop{}%
  \mkern-\thinmuskip
  \mathchoice{%
    \mint@@{#1}{#2}{#3}{#4}%
        \displaystyle\textstyle\scriptstyle
  }{%
    \mint@@{#1}{#2}{#3}{#4}%
        \textstyle\scriptstyle\scriptstyle
  }{%
    \mint@@{#1}{#2}{#3}{#4}%
        \scriptstyle\scriptscriptstyle\scriptscriptstyle
  }{%
    \mint@@{#1}{#2}{#3}{#4}%
        \scriptscriptstyle\scriptscriptstyle\scriptscriptstyle
  }%
  \mkern-\thinmuskip
  \int#1%
  \ifx\\#3\\\else_{#3}\fi
  \ifx\\#4\\\else^{#4}\fi  
}
\newcommand*{\mint@@}[7]{%
  % #1: limits
  % #2: overlay symbol
  % #3: subscript
  % #4: superscript
  % #5: math style
  % #6: math style for overlay symbol
  % #7: math style for subscript/superscript
  \begingroup
    \sbox0{$#5\int\m@th$}%
    \sbox2{$#5\int_{}\m@th$}%
    \dimen2=\wd0 %
    % => \dimen2 = width of \int
    \let\mint@limits=#1\relax
    \ifx\mint@limits\relax
      \sbox4{$#5\int_{\kern1sp}^{\kern1sp}\m@th$}%
      \ifdim\wd4>\wd2 %
        \let\mint@limits=\nolimits
      \else
        \let\mint@limits=\limits
      \fi
    \fi
    \ifx\mint@limits\displaylimits
      \ifx#5\displaystyle
        \let\mint@limits=\limits
      \fi
    \fi
    \ifx\mint@limits\limits
      \sbox0{$#7#3\m@th$}%
      \sbox2{$#7#4\m@th$}%
      \ifdim\wd0>\dimen2 %
        \dimen2=\wd0 %
      \fi
      \ifdim\wd2>\dimen2 %
        \dimen2=\wd2 %
      \fi
    \fi
    \rlap{%
      $#5%
        \vcenter{%
          \hbox to\dimen2{%
            \hss
            $#6{#2}\m@th$%
            \hss
          }%
        }%
      $%
    }%
  \endgroup
}
\newcommand{\avg}{\mint{-}}
\makeatother

% ============================================================
% Comments and drafting helpers
% ============================================================

%Comments
\newcommand{\com}[1]{\vspace{5 mm}\par \noindent
\marginpar{\textsc{Comment}} \framebox{\begin{minipage}[c]{0.95
\textwidth} \tt #1 \end{minipage}}\vspace{5 mm}\par}

% ============================================================
% Arrows
% ============================================================

\newcommand{\hooktwoheadrightarrow}{%
  \hookrightarrow\mathrel{\mspace{-15mu}}\rightarrow
}

\makeatletter
\newcommand{\xhooktwoheadrightarrow}[2][]{%
  \lhook\joinrel
  \ext@arrow 0359\rightarrowfill@ {#1}{#2}%
  \mathrel{\mspace{-15mu}}\rightarrow
}
\makeatother

\makeatletter
\newcommand*{\da@rightarrow}{\mathchar"0\hexnumber@\symAMSa 4B }
\newcommand*{\da@leftarrow}{\mathchar"0\hexnumber@\symAMSa 4C }
\newcommand*{\xdashrightarrow}[2][]{%
  \mathrel{%
    \mathpalette{\da@xarrow{#1}{#2}{}\da@rightarrow{\,}{}}{}%
  }%
}
\newcommand{\xdashleftarrow}[2][]{%
  \mathrel{%
    \mathpalette{\da@xarrow{#1}{#2}\da@leftarrow{}{}{\,}}{}%
  }%
}
\newcommand*{\da@xarrow}[7]{%
  % #1: below
  % #2: above
  % #3: arrow left
  % #4: arrow right
  % #5: space left 
  % #6: space right
  % #7: math style 
  \sbox0{$\ifx#7\scriptstyle\scriptscriptstyle\else\scriptstyle\fi#5#1#6\m@th$}%
  \sbox2{$\ifx#7\scriptstyle\scriptscriptstyle\else\scriptstyle\fi#5#2#6\m@th$}%
  \sbox4{$#7\dabar@\m@th$}%
  \dimen@=\wd0 %
  \ifdim\wd2 >\dimen@
    \dimen@=\wd2 %   
  \fi
  \count@=2 %
  \def\da@bars{\dabar@\dabar@}%
  \@whiledim\count@\wd4<\dimen@\do{%
    \advance\count@\@ne
    \expandafter\def\expandafter\da@bars\expandafter{%
      \da@bars
      \dabar@ 
    }%
  }%  
  \mathrel{#3}%
  \mathrel{%   
    \mathop{\da@bars}\limits
    \ifx\\#1\\%
    \else
      _{\copy0}%
    \fi
    \ifx\\#2\\%
    \else
      ^{\copy2}%
    \fi
  }%   
  \mathrel{#4}%
}
\makeatother

% ============================================================
% Relation and delimiter spacing
% ============================================================

\renewcommand{\geq}{\geqslant}
\renewcommand{\leq}{\leqslant}
\newcommand{\midd}{\hspace{4pt}\middle|\hspace{4pt}}

% ============================================================
% Left Right Updated Deliminators
% ============================================================

\makeatletter
\NewDocumentCommand\Left{O{}}{%
    \IfValueTF{#1}{\notblank{#1}{\mathopen\bBigg@{#1}}{\left}}{\left}}
\NewDocumentCommand\Right{O{}}{%
    \IfValueTF{#1}{\notblank{#1}{\mathopen\bBigg@{#1}}{\right}}{\right}}
\makeatother

% ============================================================
% Restrictions
% ============================================================

\newcommand{\littletaller}{\mathchoice{\vphantom{\big|}}{}{}{}}
\newcommand\restr[2]{{% we make the whole thing an ordinary symbol
  \left.\kern-\nulldelimiterspace % automatically resize the bar with \right
  #1 % the function
  \littletaller % pretend it's a little taller at normal size
  \right|_{#2} % this is the delimiter
  }}

% ============================================================
% Chapter title formatting
% ============================================================

\newcommand{\chnum}{\titleformat
{\chapter} % command
[display] % shape
{\centering} % format
{\Huge \color{black} \shadowbox{\thechapter}} % label
{-0.5em} % sep (space between the number and title)
{\LARGE \color{black} \underline} % before-code
}

\newcommand{\chunnum}{\titleformat
{\chapter} % command
[display] % shape
{} % format
{} % label
{0em} % sep
{ \begin{flushright} \begin{tabular}{r}  \Huge \color{black}
} % before-code
[
\end{tabular} \end{flushright} \normalsize
] % after-code
}



\begin{document}
\begin{center}
    {\Large Math 548 \\[0.1in]Portfolio Assignment 2}\\[.175in]
    {Name:} {\underline{Gianluca Crescenzo\hspace*{2in}}}\\[0.15in]
    \end{center}
    \vspace{4pt}
%%%%%%%%%%%%%%%%%%%%%%%%%%%%%%%%%%%%%%%%%%%%%%%%%%
\begin{problem}[F22, P4]
Let $W \subset \bfR^5$ be the space spanned by the vectors
    \begin{equation*}
    \begin{split}
        \left\{ \bmat 1 \\ -2 \\ 0 \\ 2 \\-1 \emat , \bmat -2 \\ 4 \\ -1 \\1\\2 \emat, \bmat 0\\1\\2\\-2\\1\emat\right\}.
    \end{split}
    \end{equation*}
\begin{enumerate}[label = (\alph*),itemsep=1pt,topsep=3pt]
    \item Compute the dimension of $W$.
    \item Determine the dimension of $W^\perp$, and explain how this follows immediately from (a) using a theorem.
    \item Find a basis for $W^\perp$.
\end{enumerate}
\end{problem}
    {\color{blue} \begin{proof}
        We will first put the following matrix into reduced row echelon form:
            \begin{equation*}
            \begin{split}
                \bmat 1 & -2 & 0 & 2 & -1 \\ -2 & 4 & -1 & 1 & 2 \\ 0 & 1 & 2 & -2 & 1 \emat
                & \xrightarrow{R_2 \to R_2 + 2R_1}
                \bmat 1 & -2 & 0 & 2 & -1 \\ 0 & 0 & -1 & 5 & 0 \\ 0 & 1 & 2 & -2 & 1 \emat \\[1em]
                & \xrightarrow{R_2 \leftrightarrow R_3}
                \bmat 1 & -2 & 0 & 2 & -1 \\ 0 & 1 & 2 & -2 & 1 \\ 0 & 0 & -1 & 5 & 0 \emat \\[1em]
                & \xrightarrow{R_3 \to -R_3}
                \bmat 1 & -2 & 0 & 2 & -1 \\ 0 & 1 & 2 & -2 & 1 \\ 0 & 0 & 1 & -5 & 0 \emat \\[1em]
                & \xrightarrow{R_2 \to R_2 - 2R_3}
                \bmat 1 & -2 & 0 & 2 & -1 \\ 0 & 1 & 0 & 8 & 1 \\ 0 & 0 & 1 & -5 & 0 \emat \\[1em]
                & \xrightarrow{R_1 \to R_1 + 2R_2}
                \bmat 1 & 0 & 0 & 18 & 1 \\ 0 & 1 & 0 & 8 & 1 \\ 0 & 0 & 1 & -5 & 0 \emat.
            \end{split}
            \end{equation*}

        (a) Since there are three pivots, we have $\dim W = 3$.

        (b) Since $\dim \bfR^5 = \dim W +\dim W^\perp$, we have $\dim W^\perp = 2$.

        (c) Computing the nullspace of our matrix gives:
            \begin{equation*}
            \begin{split}
                \bmat 1 & 0 & 0 & 18 & 1 \\ 0 & 1 & 0 & 8 & 1 \\ 0 & 0 & 1 & -5 & 0 \emat \bmat x_1 \\ x_2 \\ x_3 \\ x_4 \\ x_5 \emat = \bmat 0 \\ 0 \\ 0 \\0 \\0  \emat
                & = \begin{cases}
                    x_1 = 18x_3 + x_5 \\
                    x_2 = 8x_4 + x_5 \\
                    x_3 = -5x_5  
                \end{cases}
            \end{split}
            \end{equation*}
        Taking $x_4$ and $x_5$ to be free, we get:
            \begin{equation*}
            \begin{split}
                \cB_{W^\perp} = \left\{ \bmat -18 \\ -8 \\ 5 \\ 1 \\0 \emat, \bmat -1 \\ -1 \\0 \\ 0 \\ 1 \emat \right\}.
            \end{split}
            \end{equation*}
    \end{proof}}
%%%%%%%%%%%%%%%%%%%%%%%%%%%%%%%%%%%%%%%%%%%%%%%%%%
\begin{problem}[F19, P1]
    Let $P_3$ be the real vector space of all real polynomials of degree three or less. Define $L:P_3 \rightarrow P_3$ by $L(p(x)) = p(x) + p(-x)$. 
    \begin{enumerate}[label = (\alph*),itemsep=1pt,topsep=3pt]
        \item Prove that $L$ is a linear transformation.
        \item Find a basis for the null space of $L$.
        \item Compute the dimension of the image of $L$.
    \end{enumerate}
\end{problem}
    {\color{blue} \begin{proof}
        (a) Let $p(x),q(x) \in P_3$ and $\alpha \in \bfR$. Observe that:
            \begin{equation*}
            \begin{split}
                L(p(x) + \alpha q(x))
                & = p(x) + \alpha q(x) + \bigl(p(-x) + \alpha q(-x)\bigr) \\
                & = p(x) + p(-x) + \alpha(q(x) + q(-x)) \\
                & = L(p(x)) + \alpha L(q(x)).
            \end{split}
            \end{equation*}
        Thus $L$ is a linear transformation.

        (b) Let $\cB = \{1,x,x^2,x^3\}$. Then:
            \begin{equation*}
            \begin{split}
                L(1) &= 2 \\
                L(x) &= x + (-x) = 0 \\
                L(x^2) &= x^2 + (-x)^2 = 2x^2 \\
                L(x^3) &= x^3 + (-x)^3 = 0.
            \end{split}
            \end{equation*}
        This means:
            \begin{equation*}
            \begin{split}
                [L]_{\cB} = \bmat 2 & 0 & 0 & 0 \\ 0 & 0 & 0 & 0 \\ 0 & 0 & 2 & 0 \\ 0 & 0 & 0 & 0 \emat.
            \end{split}
            \end{equation*}
        Computing the nullspace of $[L]_\cB$ gives:
            \begin{equation*}
            \begin{split}
                \bmat 2 & 0 & 0 & 0 \\ 0 & 0 & 0 & 0 \\ 0 & 0 & 2 & 0 \\ 0 & 0 & 0 & 0 \emat \bmat x_1 \\ x_2 \\ x_3 \\ x_4 \emat = \bmat 0 \\ 0 \\ 0 \\ 0 \emat
                & \implies \begin{cases}
                2x_1 = 0 \\ 2x_3 = 0.
                \end{cases}
            \end{split}
            \end{equation*}
        Hence $\cB_{\ker[L]_{\cB}} = \left\{ \bmat 0 \\ 1 \\ 0 \\ 0 \emat , \bmat 0 \\ 0 \\ 0 \\ 1 \emat \right\}$.


        (c) By the rank-nullity theorem, $\dim (\Image L) = \dim (P_3) - \dim (\ker L) = 2$.
    \end{proof}}
%%%%%%%%%%%%%%%%%%%%%%%%%%%%%%%%%%%%%%%%%%%%%%%%%%
\begin{problem}[F18, P1]
    Let $T:\bfR^3 \rightarrow \bfR^3$ be the linear transformation defined by $T \left(\bmat x \\ y \\ z \emat \right) = \bmat x+y \\ 2z-x \\ y+2z \emat$.
        \begin{enumerate}[label = (\alph*),itemsep=1pt,topsep=3pt]
            \item Find a matrix that represents $T$ with respect to the standard basis for $\bfR^3$.
            \item Find a basis for the kernel of $T$.
            \item Determine the rank of $T$.
        \end{enumerate}
\end{problem}
    {\color{blue} \begin{proof}
        (a) Let $\cB$ denote the standard basis in $\bfR^3$. From the following:
            \begin{equation*}
            \begin{split}
                T \left( \bmat 1 \\ 0 \\ 0\emat \right) & = \bmat 1 \\ -1 \\ 0 \emat \\
                T \left( \bmat 0 \\ 1 \\ 0\emat \right) & = \bmat 1 \\ 0 \\ 1 \emat \\
                T \left( \bmat 0 \\ 0 \\ 1\emat \right) & = \bmat 0 \\ 2 \\ 2 \emat, \\
            \end{split}
            \end{equation*}
        we get:
            \begin{equation*}
            \begin{split}
                [T]_\cB = \bmat 1 & 1 & 0 \\ -1 & 0 & 2 \\ 0 & 1 & 2 \emat.
            \end{split}
            \end{equation*}

        (b) By row reducing $[T]_\cB$:
            \begin{equation*}
            \begin{split}
                \bmat 1 & 1 & 0 \\ -1 & 0 & 2 \\ 0 & 1 & 2 \emat
                & \xrightarrow{R_2 \to R_2 + R_1}
                \bmat 1 & 1 & 0 \\ 0 & 1 & 2 \\ 0 & 1 & 2 \emat \\[1em]
                & \xrightarrow{R_3 \to R_3 - R_2}
                \bmat 1 & 1 & 0 \\ 0 & 1 & 2 \\ 0 & 0 & 0 \emat \\[1em]
                & \xrightarrow{R_1 \to R_1 - R_2}
                \bmat 1 & 0 & -2 \\ 0 & 1 & 2 \\ 0 & 0 & 0 \emat,
            \end{split}
            \end{equation*}
            we can see:
                \begin{equation*}
                \begin{split}
                    \bmat 1 & 0 & -2 \\ 0 & 1 & 2 \\ 0 & 0 & 0 \emat \bmat x_1 \\ x_2 \\ x_3 \emat = \bmat 0 \\ 0 \\ 0 \emat 
                    & \implies \begin{cases}
                        x_1 = 2x_3 \\ x_2 = -2x_3.
                    \end{cases}
                \end{split}
                \end{equation*}
            Taking $x_3$ to be free, a basis for $\ker T$ is $\left\{ \bmat 2 \\ -2 \\ 1 \emat \right\}$.

        (c) The rank of a matrix is equal to the dimension of its image. By the rank-nullity theory, $\dim(\Image T) = \dim(\bfR^3) - \dim (\ker T) = 2$.
    \end{proof}}
%%%%%%%%%%%%%%%%%%%%%%%%%%%%%%%%%%%%%%%%%%%%%%%%%%
\begin{problem}[S17, P1]
    Let $V = \{a_0 + a_1 \sqrt[3]{2} + a_2 \sqrt[3]{4} \mid a_1,a_2,a_2 \in \bfQ\} \subseteq \bfR$. This is a vector space over $\bfQ$.
        \begin{enumerate}[label = (\alph*),itemsep=1pt,topsep=3pt]
            \item Verify $V$ is closed under product (using the usual product operation in $\bfR$).
            \item Let $T:V \rightarrow V$ be the linear transformation defined by $T(v) = (\sqrt[3]{2} + \sqrt[3]{4})v$. Find the matrix that represents $T$ with respect to the basis $\{1, \sqrt[3]{2}, \sqrt[3]{4}\}$ for $V$.
            \item Determine the characteristic polynomial for $T$.
        \end{enumerate}
\end{problem}
    {\color{blue} \begin{proof}
    (a) Let $a_0 + a_1 \sqrt[3]{2} + a_2 \sqrt[3]{4}$ and $b_0 + b_1 \sqrt[3]{2} + b_2 \sqrt[3]{4}$ be elements of $V$. We have:
        \begin{equation*}
        \begin{aligned}
            &(a_0 + a_1 \sqrt[3]{2} + a_2 \sqrt[3]{4})(b_0 + b_1 \sqrt[3]{2} + b_2 \sqrt[3]{4}) \\
            &\hspace{2.6em}= (a_0b_0 + 2a_1b_2 + 2a_2b_1)
            + \sqrt[3]{2}(a_0b_1 + a_1b_0 + 2a_2b_2)
            + \sqrt[3]{4}(a_0b_2 + a_1b_1 + a_2b_0).
        \end{aligned}
        \end{equation*}
    Thus $V$ is closed under product.

    (b) Let $\cB = \{1,\sqrt[3]{2},\sqrt[3]{4}\}$. Then:
        \begin{equation*}
        \begin{split}
            T(1) &= \sqrt[3]{2} + \sqrt[3]{4} \\
            T(\sqrt[3]{2}) &= 2 + \sqrt[3]{4} \\
            T(\sqrt[3]{4}) & = 2 + 2\sqrt[3]{2}.
        \end{split}
        \end{equation*}
    This means:
        \begin{equation*}
        \begin{split}
            [T]_{\cB} = \bmat 0 & 2 & 2 \\ 1 & 0 & 2 \\ 1 & 1 & 0 \emat.
        \end{split}
        \end{equation*}

    (c) The characteristic polynomial of $[T]_\cB$ is:
        \begin{equation*}
        \begin{split}
            \det\bigl([T]_\cB - \lambda I\bigr)
            & = \det \bmat -\lambda & 2 & 2 \\ 1 & -\lambda & 2 \\ 1 & 1 & -\lambda \emat \\
            & = -\lambda \det \bmat -\lambda & 2\\ 1 & -\lambda \emat - 2 \det \bmat 1 & 2 \\ 1 & -\lambda \emat + 2 \det \bmat 1 & -\lambda \\ 1 & 1 \emat \\
            & = -\lambda (\lambda^2 - 2) - 2(-\lambda -2) + 2(1 + \lambda) \\
            & = -\lambda^3 + 6\lambda +6. \qedhere
        \end{split}
        \end{equation*}
    \end{proof}}
%%%%%%%%%%%%%%%%%%%%%%%%%%%%%%%%%%%%%%%%%%%%%%%%%%
\begin{problem}[F19, P2]
    Let $T:\bfR^3 \rightarrow \bfR^3$ be the linear transformation that rotates counterclockwise around the $z$-axis by $\frac{2\pi}{3}$. 
    \begin{enumerate}[label = (\alph*),itemsep=1pt,topsep=3pt]
        \item Write the matrix for $T$ with respect to the standard basis $\left\{ \bmat 1 \\ 0 \\ 0 \emat, \bmat 0 \\ 1 \\ 0 \emat,\bmat 0 \\ 0 \\ 1 \emat \right\}$.
        \item Write the matrix for $T$ with respect to the basis $\left\{ \bmat \frac{\sqrt{3}}{2} \\ -\frac{1}{2} \\ 0 \emat, \bmat 0 \\ 1 \\ 0 \emat,\bmat 0 \\ 0 \\ 1 \emat \right\}$.
        \item Determine all (complex) eigenvalues of $T$.
        \item Is $T$ diagonalizable over $\bfC$? Justify your answer.
    \end{enumerate}
\end{problem}
    {\color{blue} \begin{proof}
        (a) Let $\cE$ denote the standard basis. We have:
            \begin{equation*}
            \begin{split}
                T \left( \bmat 1 \\ 0 \\ 0 \emat \right) & = \bmat -\frac{1}{2} \\ \frac{\sqrt{3}}{2} \\ 0 \emat, \\
                T \left( \bmat 0 \\ 1 \\ 0 \emat \right) & =\bmat -\frac{\sqrt{3}}{2} \\ -\frac{1}{2} \\ 0 \emat, \\
                T \left( \bmat 0 \\ 0 \\ 1 \emat \right) & = \bmat 0 \\ 0 \\ 1 \emat. 
            \end{split}
            \end{equation*}
        Hence:
            \begin{equation*}
            \begin{split}
                [T]_\cE = \bmat 
                -\frac{1}{2} & -\frac{\sqrt{3}}{2} & 0 \\ 
                \frac{\sqrt{3}}{2} & -\frac{1}{2} & 0 \\ 
                0 & 0 & 1 
                \emat.
            \end{split}
            \end{equation*}

        (b) Let $\cB = \left\{ \bmat \frac{\sqrt{3}}{2} \\ -\frac{1}{2} \\ 0 \emat, \bmat 0 \\ 1 \\ 0 \emat, \bmat 0 \\ 0 \\ 1 \emat \right\}$. We have:
            \begin{equation*}
            \begin{split}
                T \left( \bmat \frac{\sqrt{3}}{2} \\ -\frac{1}{2} \\ 0 \emat \right)
                & = \bmat 0 \\ 1 \\ 0 \emat, \\
                T \left( \bmat 0 \\ 1 \\ 0 \emat \right)
                & = \bmat -\frac{\sqrt{3}}{2} \\ -\frac{1}{2} \\ 0 \emat, \\
                T \left( \bmat 0 \\ 0 \\ 1 \emat \right)
                & = \bmat 0 \\ 0 \\ 1 \emat.
            \end{split}
            \end{equation*}
        Hence:
            \begin{equation*}
            \begin{split}
                [T]_\cB = 
                \bmat 
                0 & -1 & 0 \\ 
                1 & 1 & 0 \\ 
                0 & 0 & 1 
                \emat.
            \end{split}
            \end{equation*}

        (c) We must first find the characteristic polynomial of $[T]_\cB$:
            \begin{equation*}
            \begin{split}
                \det \bigl([T]_\cB - \lambda I\bigr)
                & = \det \bmat -\lambda & -1 & 0 \\ 1 & 1-\lambda & 0 \\ 0 & 0 & 1-\lambda \emat \\
                & = (1-\lambda)\det \bmat -\lambda & -1 \\ 1 & 1-\lambda \emat \\
                & = (1-\lambda)(\lambda^2 - \lambda + 1). \\
            \end{split}
            \end{equation*}
        Hence $\lambda = 1$, $\lambda = \frac{1}{2} + \frac{\sqrt{3}}{2}i$, and $\lambda = \frac{1}{2} - \frac{\sqrt{3}}{2}i$ are the eigenvalues of $[T]_\cB$.

        (d) Since $T$ has three distinct eigenvalues, and since $\dim(\bfR^3) = 3$, $T$ is diagonalizable over $\bfC$.

    \end{proof}}
%%%%%%%%%%%%%%%%%%%%%%%%%%%%%%%%%%%%%%%%%%%%%%%%%%
\end{document}